which shows that, for @@M0@@, @@M1@@ maps to the @@M2@@-th symmetric function of the @@M3@@, and consequently that @@M4@@ is injective and identifies @@M5@@ with the ring of symmetric polynomials in the @@M6@@.

We then have @@M7@@, and @@M8@@. The @@M9@@ are polynomials in the @@M10@@; if, for some index @@M11@@, we set @@M12@@, we see that @@M13@@ becomes @@M14@@, hence that the @@M15@@ vanish for @@M16@@, and therefore that they contain @@M17@@ as a factor. Thus they contain the product of the @@M18@@ as a factor, i.e. they are divisible by @@M19@@ in @@M20@@, and hence also in @@M21@@, since the quotient of two symmetric polynomials is evidently symmetric.

\sgaStableTarget{sga6.e05.def.5-3}5.3. Since @@M22@@ is an integral domain, we define an element @@M23@@ unambiguously by putting (@@M24@@):
\sgaStableTarget{sga6.e05.eq.5-3-1}@@M25@@

The element @@M26@@ is therefore a polynomial in the @@M27@@, @@M28@@, isobaric of weight @@M29@@ with respect to the @@M30@@, since the same is true of @@M31@@ by 2.3.

Let @@M32@@ be an integer @@M33@@, and let @@M34@@ be the @@M35@@-@@M36@@-algebra generated by the generators @@M37@@ and @@M38@@ subject to the relations @@M39@@ for @@M40@@; by Proposition 4.7, there exists exactly one @@M41@@-homomorphism @@M42@@, sending @@M43@@ to @@M44@@ and @@M45@@ to @@M46@@. Definition (5.3.1) implies that @@M47@@ for every @@M48@@.

If @@M49@@ is a @@M50@@-ring and @@M51@@ is an element of @@M52@@ such that @@M53@@ for @@M54@@, we can, for every @@M55@@, define an element @@M56@@ by substitution, or by using the unique @@M57@@-homomorphism from @@M58@@ to
